% SPDX-License-Identifier: 0BSD
\section{\texorpdfstring{Worked \(d\bar d\to Zgg\) Example}{Worked d dbar -> Z g g Example}}

The process used here is
\[
  d(p_1)\,\bar d(p_2)\to Z(p_3)\,g(p_4)\,g(p_5).
\]
The report source labels are the process-order labels \(1,\ldots,5\). Labels
\(1\) and \(2\) are the incoming quark and antiquark, label \(3\) is the
colour-singlet \(Z\), and labels \(4,5\) are the two final-state gluons. The
leading-colour selected-flow example keeps the historical colour traversal
\[
  (2,5,4,1),
\]
meaning that the open colour line starts at the antiquark, crosses the two
gluons in the order \(g_5,g_4\), and closes on the quark. In the public
\PAC{} selection object this is recorded as the source-label order
\((2,5,4,1,3)\): the final \(3\) keeps the colour-singlet source label attached
to the benchmark metadata, while colour-sector matching ignores it when
building the coloured word.

The same physical process is used in two LC workloads. The selected-flow,
helicity-sum workload generates only the LC sector compatible with
\((2,5,4,1)\) and sums over all source helicities. The all-flows,
fixed-helicity workload keeps all LC colour flows, such as the sibling order
\((2,4,5,1)\), and fixes one helicity assignment that is present in the
generated DAG. Both workloads are compared to the same original Fortran
\AC{} generation. The all-flow fixed-helicity reference does not require a
second \AC{} generation step; only its runtime observable differs.

\begin{center}
\small
\begin{tabular}{@{}L{1.15in}L{5.65in}@{}}
\toprule
\textbf{object} & \textbf{report convention} \\
\midrule
source labels & \(1=d_{\rm in},\ 2=\bar d_{\rm in},\ 3=Z,\ 4=g,\ 5=g\) \\
selected coloured word & \((2,5,4,1)\), with the \(Z\) omitted from colour matching \\
public selected order & \((2,5,4,1,3)\), stored in the report provenance \\
selected LC sectors & the physical LC sector and any required internal sibling sectors for the same public flow \\
runtime observable & summed matrix element from \texttt{Runtime.evaluate()} at fixed phase-space points \\
\bottomrule
\end{tabular}
\end{center}

The old \AC{} reference and the \PAC{} measurements use aligned numerical
input parameters and the same deterministic point set. Any disagreement in the
selected value, the all-flow fixed-helicity value, or the resolved-sum identity
is therefore a validation failure, not a timing fluctuation. Timing entries are
accepted only after this numerical comparison has passed.

